%% kapa: exam=1 round=32 year=2021 period=1 form=A subject=real_estate subject_ko=부동산학원론 seq=32 number=112
%% source: Q-Net 감정평가사 기출문제 (cst00309) artlSeq=5208397 file=extracted/1008/5208397/A형/2021년도 제32회 감정평가사 1차 1교시 A형.pdf
%% conversion: kordoc 4.24.51 -> tools/gen1.py (auto); review=proofread-sonnet; flags=none
\begin{question}{112}
부동산 권리분석에 관련된 설명으로 옳지 않은 것은?
\begin{choices}
  \item 부동산 권리관계를 실질적으로 조사·확인·판단하여 일련의 부동산활동을 안전하게 하려는 것이다.
  \item 대상부동산의 권리관계를 조사·확인하기 위한 판독 내용에는 권리의 하자나 거래규제의 확인·판단이 포함된다.
  \item 매수인이 대상부동산을 매수하기 전에 소유권이전을 저해하는 사항이 있는지 여부를 확인하기 위하여 공부(公簿) 등을 조사하는 일도 포함된다.
  \item 우리나라 등기는 관련 법률에 다른 규정이 있는 경우를 제외하고는 당사자의 신청 또는 관공서의 촉탁에 따라 행하는 신청주의 원칙을 적용한다.
  \item 부동산 권리분석을 행하는 주체가 분석대상권리의 주요한 사항을 직접 확인해야 한다는 증거주의의 원칙은 권리분석활동을 하는 데 지켜야 할 이념이다.
\end{choices}
\end{question}
\answer{5}
